\documentclass[11pt]{article}

\usepackage[margin=2.5cm]{geometry}
\usepackage{graphicx}
\usepackage{amsmath}
\usepackage{amssymb}
\usepackage{booktabs}
\usepackage{hyperref}
\usepackage{caption}
\usepackage{xcolor}
\usepackage{natbib}
\usepackage{float}

\title{QML power-spectrum estimation on the TR1 South field:\\
smaller errors, clean low-$\ell$ $B$-modes, and what it took}
\author{Weak Lensing QML pipeline}
\date{\today}

\begin{document}
\maketitle

\begin{abstract}
We apply the conjugate-gradient Quadratic Maximum Likelihood (QML) estimator of \citet{Maraio2022} to the real
TR1 South shear catalogue, with all six tomographic bins combined (21.9 galaxies\,arcmin$^{-2}$ effective), on a
footprint of only $f_{\rm sky}=0.90\%$ at $N_{\rm side}=64$. We compare it with two independent pseudo-$C_\ell$
estimates: NaMaster on the same map, and heracles run directly on the catalogue. The $E$-mode bandpowers of all three
estimators agree with each other and with the fiducial model ($\chi^2=7.7$ for 7 bands). QML's errors are smaller
than pseudo-$C_\ell$'s in every band, by $3\times$ for $EE$ and by $28\times$ ($8\times$) for $BB$ at $\ell\approx20$
($40$). With these much tighter errors, QML's lowest-$\ell$ $B$-modes are consistent with zero, while the
pseudo-$C_\ell$ estimators show mask-induced $E\to B$ leakage there. Two things made this work. First, a 50\% pixel-coverage cut:
partly covered edge pixels have up to $10^4\times$ the noise QML assumes, which made the QML estimate fail completely
($\chi^2\sim10^4$), while heracles, which weights pixels by their coverage, barely changed. Second, bandpowers of width
$\Delta\ell=20$ formed before inverting the Fisher matrix. Re-analysing our GLASS $N_{\rm side}=64$ mock grid
the same way shows that the large-footprint requirement reported for those mocks comes from per-$\ell$ estimation:
with bandpowers, even the narrowest (15\%) footprint gives clean $B$-modes. The natural next step is a per-pixel noise
model in QML, which would remove the need for the coverage cut.
\end{abstract}

\tableofcontents

\section{Data and method}
\label{sec:method}

\paragraph{Data.} The TR1 South weak-lensing catalogue (31.0 million galaxies, MetaCal shapes with the mean response
already applied per tomographic bin) and its effective-coverage map. \textbf{All six tomographic bins are combined}:
a single bin has only 3.3--4.1 galaxies\,arcmin$^{-2}$, so the combined sample is far less noisy. After the coverage
cut of \S\ref{sec:cut}, the effective density is $n_{\rm eff}=21.9$\,arcmin$^{-2}$ with a per-component shape
dispersion $\sigma_\epsilon=0.334$. In shot-noise terms ($\sigma_\epsilon^2/n$) that equals a sample with
$\sigma_\epsilon=0.3$ at $n\approx17.7$\,arcmin$^{-2}$, close to the densest case of our GLASS mocks.

\paragraph{Maps and mask.} Weighted mean shear per HEALPix pixel at $N_{\rm side}=64$ (55$'$ pixels). The mask
keeps pixels at least 50\% covered: 442 pixels, $f_{\rm sky}=0.90\%$, a patch about $20^\circ$ across.

\paragraph{Estimators.} All three use the same mask and the same bandpowers:
\begin{itemize}
\item \textbf{QML:} the C++ conjugate-gradient implementation, with a Monte Carlo Fisher matrix (10 Gaussian maps per
$\ell$) up to $\ell_{\max}=150$ and the fiducial cosmology of \citet{Maraio2022}. The fiducial $C^{EE}_\ell$ comes from
CCL with $n(z)$ taken from the photometric redshifts. The shot noise is a single per-pixel variance
$\sigma_\epsilon^2/\bar n_{\rm pix}$.
\item \textbf{NaMaster pseudo-$C_\ell$} on the same map, with its Gaussian covariance.
\item \textbf{heracles}, run directly on the catalogue with its own mixing matrices. It keeps the shot noise
(\texttt{debias=False}), so that it is on the same footing as QML and NaMaster, which both estimate signal plus
noise. The noise is subtracted from all three in the plots.
\end{itemize}

\paragraph{Bandpowers.} Seven bands of width $\Delta\ell=20$ from $\ell=10$, the last being 130--150. For QML the
per-$\ell$ Fisher matrix $\mathbf F$ and quadratic statistic $\mathbf y$ are binned \emph{before} inverting:
$\mathbf F_b=\mathbf W\mathbf F\mathbf W^{T}$, $\hat{\mathbf c}_b=\mathbf F_b^{-1}\mathbf W\mathbf y$, and the bandpower
covariance is ${\rm cov}(\hat{\mathbf c}_b)=\mathbf F_b^{-1}$, with $W_{b\ell}=1$ for $\ell$ in band $b$.

\paragraph{B-modes as a null test.} Weak lensing produces essentially no $B$-modes, so the measured $BB$ (after
removing the noise) should be consistent with zero. A significant $B$-mode signal points to leakage from the
estimator or to a systematic in the data. We judge this by comparing the $BB$ bandpowers with zero relative to their
errors ($\chi^2$ and its probability to exceed, PTE).

The full analysis is in \texttt{TR1/TR1\_south\_power\_spectra\_and\_covariance\_estimation.ipynb}; the figures here
are made by \texttt{TR1/report\_south/make\_figures.py}.

\section{Main results}
\label{sec:results}

\subsection{QML errors are smaller than pseudo-$C_\ell$'s in every band}

Figure~\ref{fig:precision} shows the ratio of the pseudo-$C_\ell$ (NaMaster) to QML bandpower errors. QML is more
precise everywhere. The gain is largest on the largest scales: $28\times$ for $BB$ and $3\times$ for $EE$ at
$\ell\approx20$. It falls to about $1\times$ ($BB$) and $1.2\times$ ($EE$) by $\ell\approx100$. This is the pattern
\citet{Maraio2022} predict: mask-induced mode coupling and $E/B$ mixing matter most on large scales, where the signal
dominates the noise.

\begin{figure}[H]
\centering
\includegraphics[width=0.72\textwidth]{figures/fig1_precision.png}
\caption{Ratio of pseudo-$C_\ell$ to QML bandpower errors for TR1 South ($EE$ and $BB$). Above 1, QML is more precise.}
\label{fig:precision}
\end{figure}

\subsection{$B$-modes}

Figure~\ref{fig:bmodes} shows the $BB$ bandpowers. At $\ell\lesssim50$ QML's $B$-modes are consistent with zero
($+1.1\sigma$ and $-0.6\sigma$) with errors 8--28$\times$ smaller than pseudo-$C_\ell$'s. heracles' lowest band sits
$7.9\sigma_{\rm QML}$ above zero: leakage of the strong large-scale $E$ power into $B$ through the mask. NaMaster has
the same excess, hidden by its much larger error bar. Being able to measure large-scale $B$-modes this precisely,
without leakage, is the main advantage of QML on a small patch.

The $B$-mode test is not clean everywhere:
\begin{itemize}
\item \textbf{$\ell=50$--69:} $+3.1\sigma$.
\item \textbf{$\ell\geq110$:} $+3.4\sigma$ and $+3.9\sigma$. This excess is shared by heracles and NaMaster, so it is
in the data or its pixelisation rather than specific to QML. The most likely cause is aliasing of small-scale power
into these multipoles at $N_{\rm side}=64$ (the band edge is near $2N_{\rm side}$), but a real systematic can't be
excluded yet (\S\ref{sec:next}).
\end{itemize}
Over all seven bands, QML's $\chi^2$ against zero is $50.4$; over $\ell=10$--109 it is $15.0$ for 5 bands (PTE 0.01),
driven by the 50--69 band.

\begin{figure}[H]
\centering
\includegraphics[width=0.8\textwidth]{figures/fig2_bmodes.png}
\caption{TR1 South $BB$ bandpowers (noise removed): QML (with errors), NaMaster pseudo-$C_\ell$ (with errors, which
extend beyond the panel at $\ell\lesssim40$), and heracles on the catalogue.}
\label{fig:bmodes}
\end{figure}

\subsection{$E$-modes}

All three estimators agree with each other and with the fiducial theory (Figure~\ref{fig:emodes}). For QML,
$\chi^2=7.7$ for 7 bands (PTE 0.36). In the lowest band QML sits on the theory curve, while both pseudo-$C_\ell$
estimators are low, consistent with $E$ power leaking into $B$ there.

\begin{figure}[H]
\centering
\includegraphics[width=0.8\textwidth]{figures/fig3_emodes.png}
\caption{TR1 South $EE$ bandpowers (noise removed) from the three estimators, with the fiducial signal.}
\label{fig:emodes}
\end{figure}

\begin{table}[H]
\centering
\begin{tabular}{lccccccc}
\toprule
band centre $\ell$ & 20 & 40 & 60 & 80 & 100 & 120 & 140 \\
\midrule
$\sigma^{\rm PCl}/\sigma^{\rm QML}$, $EE$ & 3.0 & 1.6 & 1.3 & 1.2 & 1.2 & 1.2 & 1.3 \\
$\sigma^{\rm PCl}/\sigma^{\rm QML}$, $BB$ & 27.8 & 7.8 & 2.6 & 1.4 & 1.1 & 1.0 & 1.0 \\
QML $BB/\sigma^{\rm QML}$ & 1.1 & $-0.6$ & 3.1 & 1.3 & $-0.7$ & 3.4 & 3.9 \\
heracles $BB/\sigma^{\rm QML}$ & 7.9 & 2.4 & 2.7 & 2.5 & 0.2 & 6.6 & 7.8 \\
\bottomrule
\end{tabular}
\caption{Bandpower precision and $B$-mode null test for TR1 South (noise removed from $BB$).}
\label{tab:results}
\end{table}

\section{How the result was obtained}
\label{sec:how}

A first run, on the same data with the TR1 north notebook's settings, gave QML bandpowers 10--100$\times$ too large.
Four changes fixed it. The first two account for almost all of the improvement.

\subsection{Removing partly covered pixels}
\label{sec:cut}

At $N_{\rm side}=64$ many pixels on the edge of the survey, and along internal gaps, are only partly covered
(Figure~\ref{fig:coverage}). They hold few galaxies, so their noise is far higher than average. QML, however, assumes
a \emph{single} noise variance for every pixel. Figure~\ref{fig:pixelnoise} shows the real per-pixel noise variance
($\sigma_\epsilon^2\sum w^2/(\sum w)^2$) divided by that assumed value:
\begin{itemize}
\item \textbf{Pixels covered less than 50\%:} median $4\times$, mean $239\times$, worst pixel $13{,}300\times$.
\item \textbf{Kept pixels:} median 0.98, mean 1.02, maximum 1.7.
\end{itemize}
The cut removes 96 of 538 pixels (18\%). These are only 5.2\% of the covered area, and 5.1\% of the galaxies.

\begin{figure}[H]
\centering
\includegraphics[width=\textwidth]{figures/fig4_coverage_cut.png}
\caption{Left: TR1 South pixels at $N_{\rm side}=64$ kept (blue) and removed (orange) by the 50\% coverage cut: the
survey edge and internal gaps. Right: covered fraction of every pixel touching the survey, with the cut.}
\label{fig:coverage}
\end{figure}

\begin{figure}[H]
\centering
\includegraphics[width=0.72\textwidth]{figures/fig5_pixel_noise.png}
\caption{Real per-pixel noise variance divided by the single value QML assumes, against covered fraction.}
\label{fig:pixelnoise}
\end{figure}

\subsection{QML depends on the cut; heracles does not}

Figure~\ref{fig:cut} compares both estimators with and without the cut.
\begin{itemize}
\item \textbf{QML without the cut} is badly wrong: $EE$ is 3--10$\times$ above theory and $BB$ is 10--100$\times$
above zero, with $\chi^2=13{,}407$ ($EE$) and $32{,}228$ ($BB$) for 7 bands.
\item \textbf{QML with the cut:} $\chi^2=7.7$ and $50.4$.
\item \textbf{heracles barely changes.} It works from the catalogue's galaxy weights and the fractional coverage map,
so partly covered pixels are automatically down-weighted.
\end{itemize}
So the edge pixels hold no anomalous signal. The failure was entirely QML's uniform-noise assumption, which is what
the per-pixel noise model of \S\ref{sec:next} would fix properly.

\begin{figure}[H]
\centering
\includegraphics[width=\textwidth]{figures/fig6_cut_effect.png}
\caption{QML (circles) and heracles (squares) with (filled) and without (open) the 50\% coverage cut. Noise removed;
symmetric-log axis. Left: $EE$ with the fiducial signal; right: $BB$ with zero.}
\label{fig:cut}
\end{figure}

\subsection{Bandpowers at least as wide as the patch can resolve}

A $20^\circ$ patch can't separate multipoles closer than $\Delta\ell\approx2\pi/\theta\approx18$. The north
notebook's binning (equal numbers of modes per band) made the high-$\ell$ bands only 5--7 multipoles wide.
Neighbouring QML bandpowers were then almost perfectly anticorrelated (correlation down to $-0.92$) and alternated in
sign by $\pm5\sigma$. Uniform bands of width 20 reduce the neighbouring correlations to about $-0.2$.

\subsection{Noise accounting}

The shot-noise level and density are computed from the galaxies in kept pixels only. Counting the removed galaxies
over the kept area had overstated the density (23.1 instead of 21.9\,arcmin$^{-2}$) and understated QML's noise
variance by 5\%.

\section{Why such a small footprint works: comparison with the GLASS mocks}
\label{sec:glass}

Our GLASS $N_{\rm side}=64$ mock grid (log-normal catalogues; densities $n=1,10,20$\,arcmin$^{-2}$ with
$\sigma_\epsilon=0.3$; footprints with $f_{\rm sky}=15$, 31 and 56\%) suggested that QML needs both a large footprint
and a high density: at $f_{\rm sky}=15\%$ the $B$-modes were leakage-dominated at every density. TR1 South has
$f_{\rm sky}=0.9\%$, 17 times smaller than the smallest GLASS footprint.

That GLASS assessment used \emph{per-$\ell$} QML estimates. Each multipole is then estimated by inverting a nearly
singular, Monte-Carlo-estimated Fisher matrix, which is unstable when the sky is heavily masked. We re-analysed the
same GLASS Fisher matrices and $y_\ell$ the TR1 way, with $\Delta\ell=20$ bandpowers formed before inversion
(Table~\ref{tab:glass}, Figure~\ref{fig:glass}). \textbf{Every GLASS case then gives a clean $B$-mode null test,
including the 15\% footprint at all three densities.}

\begin{table}[H]
\centering
\small
\begin{tabular}{llccc}
\toprule
footprint & $n$\,[arcmin$^{-2}$] & per-$\ell$ $|BB|$/noise & bandpower $|BB|$/noise & bandpower PTE \\
\midrule
(55,125)$^\circ$, $f_{\rm sky}=15\%$ & 1 / 10 / 20 & 75.1 / 30.4 / 30.8 & 0.95 / 1.11 / 1.03 & 0.18 / 0.13 / 0.68 \\
(65,115)$^\circ$, $f_{\rm sky}=31\%$ & 1 / 10 / 20 & 1.0 / 1.5 / 1.2 & 0.89 / 1.07 / 1.07 & 0.03 / 0.47 / 0.56 \\
(75,105)$^\circ$, $f_{\rm sky}=56\%$ & 1 / 10 / 20 & 0.9 / 1.0 / 1.1 & 0.98 / 0.98 / 1.04 & 0.11 / 0.43 / 0.46 \\
\bottomrule
\end{tabular}
\caption{GLASS $N_{\rm side}=64$ grid, QML $B$-mode null test. Per-$\ell$: the metric of the GLASS report (mean over
$\ell\leq64$). Bandpowers: $\Delta\ell=20$ bands covering $\ell=10$--69, with the $\chi^2$ PTE against pure noise.}
\label{tab:glass}
\end{table}

\begin{figure}[H]
\centering
\includegraphics[width=\textwidth]{figures/fig7_glass_plane.png}
\caption{The GLASS grid in the sky-fraction / density plane (density scaled to $\sigma_\epsilon=0.3$), coloured by
whether the QML $B$-mode null test is clean. Left: per-$\ell$ estimates, as in the GLASS report. Right: $\Delta\ell=20$
bandpowers. The filled star is TR1 South with all bins combined; open stars are its individual tomographic bins.}
\label{fig:glass}
\end{figure}

So TR1 South works because of three things together:
\begin{enumerate}
\item it is estimated as \textbf{bandpowers}, which removes the per-$\ell$ instability behind the footprint
requirement;
\item it is \textbf{dense}: combining all bins puts it close to the best GLASS density;
\item after the coverage cut its \textbf{noise really is uniform}, as QML assumes.
\end{enumerate}
Two caveats apply. The GLASS bandpower test covers only $\ell=10$--69 (three bands), and footprints of at least 15\%.
Going down to $f_{\rm sky}=0.9\%$ is supported by the TR1 results themselves (the $E$-mode agreement and the
consistency between estimators) rather than by mocks at that size. Mocks matching the TR1 South footprint would close
this gap.

\section{Caveats}
\label{sec:caveats}
\begin{itemize}
\item \textbf{All tomographic bins are combined.} This is a single-sample measurement, not a tomographic analysis.
\item \textbf{The QML Fisher matrix is a Monte Carlo estimate from 10 maps per $\ell$} (the paper used 25). It is
seeded from the clock, so runs aren't exactly reproducible. With bandpowers the matrix is positive definite, but its
scatter enters the errors at the 10--20\% level at low $\ell$.
\item \textbf{The kept pixels' noise still varies} by 0.85--1.6$\times$ about the assumed value, increasing towards
the cut.
\item \textbf{The $B$-mode excess at $\ell\geq110$} is common to all three estimators and unexplained.
\item \textbf{The MetaCal selection response} can't be read from the catalogue. It would change the shear amplitude
(the $EE$ normalisation), not the noise or $B$-mode results.
\item The fiducial cosmology is fixed; $\ell<10$ is not used.
\end{itemize}

\section{Next steps}
\label{sec:next}
\begin{enumerate}
\item \textbf{Variable (per-pixel) noise in QML.} Replace the single \texttt{noise\_var} with a map of each pixel's
noise variance, $\sigma_\epsilon^2\sum w^2/(\sum w)^2$. In the C++ this is contained: the noise term of the
covariance product (\texttt{noise\_var * x} in \texttt{cpp/detail/MatrixReplacement.h}) becomes an element-wise
product with a per-pixel vector, read in alongside the mask. QML could then keep every pixel, weighting each
correctly as heracles does. This removes the coverage cut and the remaining 0.85--1.6$\times$ variation, and
recovers the 5\% of galaxies the cut discards. \citet{Maraio2022} identify exactly this as a strength of pixel-based
QML over pseudo-$C_\ell$.
\item \textbf{Per-tomographic-bin analysis.} Each bin has 3.3--4.1 galaxies\,arcmin$^{-2}$, equivalent to
$n\approx2.3$--3.6 at $\sigma_\epsilon=0.3$. The re-binned GLASS results pass the null test even at $n=1$, so a
per-bin analysis looks feasible, at larger relative errors.
\item \textbf{The high-$\ell$ $B$-mode excess:} repeat at $N_{\rm side}=128$. If the excess moves with the pixel
scale it is aliasing; if it stays, it is in the data.
\item \textbf{Fisher accuracy:} more Monte Carlo maps (25--50) with fixed seeds, or, for a mask this small (884 data
values), the exact Fisher matrix computed directly, removing the Monte Carlo noise altogether.
\item \textbf{TR1 north:} rerun with the same pipeline (coverage cut, $\Delta\ell\approx25$--30 for its smaller patch).
\end{enumerate}

\begin{thebibliography}{9}
\bibitem[Maraio, Hall \& Taylor(2022)]{Maraio2022}
A.~Maraio, A.~Hall, A.~N.~Taylor,
\emph{Testing Quadratic Maximum Likelihood estimators for forthcoming Stage-IV weak lensing surveys},
MNRAS \textbf{000}, 1--14 (2022), arXiv:2207.10412.
\end{thebibliography}

\end{document}
